%% Journal of Econometrics submission, Elsevier elsarticle (Harvard/authoryear).
%% Template: elsarticle-template-harv.tex (elsarticle 3.x).
\documentclass[preprint,12pt,authoryear,times]{elsarticle}

\usepackage{amssymb,mathtools,amsthm,booktabs,float,siunitx}
\usepackage[hidelinks]{hyperref}
\mathtoolsset{showonlyrefs}

\journal{Journal of Econometrics}

\newtheorem{theorem}{Theorem}
\newtheorem{remark}{Remark}

\newcommand{\E}{\mathbb{E}}
\newcommand{\Q}{\mathbb{Q}}
\newcommand{\R}{\mathbb{R}}

\begin{document}

\begin{frontmatter}

\title{\texorpdfstring{Thricing Is All You Need\\[0.35em]
{\small Well, and a proper bandwidth: risk neutral densities and accurate prices from the same kernel}}{Thricing Is All You Need}}

\author[ifz]{Simon A.\ Broda\corref{cor1}}
\ead{simon.broda@hslu.ch}
\affiliation[ifz]{organization={Institute of Financial Services IFZ, Lucerne University of Applied Sciences},
            addressline={Suurstoffi 1},
            city={Rotkreuz},
            postcode={6343},
            country={Switzerland}}
\cortext[cor1]{Corresponding author.}

\begin{abstract}
Kernel estimators of the risk-neutral density differentiate a smooth of the call twice. The bandwidth that prices the quotes leaves many peaks. The bandwidth that removes the peaks moves the price. This paper uses one Gaussian kernel, at the second-derivative rate, on a cubic spline of the complementary distribution of a single maturity, and thrices that kernel. The same fit is a unimodal density of unit mass and a call interpolant. The convolution has a closed form, and the estimator is in-fill consistent. On four listed chains the out-of-the-money error is below the lognormal mixture and with the pricing-bandwidth kernels.
\end{abstract}

\begin{keyword}
risk-neutral density \sep option prices \sep kernel estimation \sep in-fill consistency \sep state-price density
\MSC[2010] 62G07 \sep 62G20 \sep 62P05 \sep 91G20
\JEL C14 \sep C58 \sep G13
\end{keyword}

\end{frontmatter}

\section{Introduction}
\label{sec:intro}

\citet{breeden1978} identified the risk-neutral density of the underlying at a fixed expiry with the second strike derivative of the corresponding European call, scaled by a discount factor. The identity assumes a twice differentiable price curve. A market supplies a finite ordered list of strikes and a quote at each knot. Numerical second differentiation of that list is ill-posed: interpolants oscillate, change sign, and depend on the interpolating scheme. The literature responds with parametric mixtures and Gram--Charlier expansions, implied-volatility interpolants, local polynomials of the call, and \citet{bondarenko2003}'s positive convolution approximation, which forces the density to be a mixture of a positive kernel with a discrete probability fitted to prices. Section~\ref{sec:lit} reviews these approaches.

Kernel smoothers of the call are usually run at two bandwidths. The width that minimises a pricing error tracks the quotes, and the second derivative has many peaks. The width that balances the variance of a second derivative, of order \(n^{-1/9}\), is the one the surveys widen further, and the price moves with it \citep{aitsahalia1998,aitsahalia2003,grith2012,jackwerth2004,figlewski2018}.

This paper uses one bandwidth for both objects, that rate with a short constant, on a single maturity. The calls are rescaled to a complementary distribution, a cubic spline interpolates that function, and a Gaussian kernel smooths the spline. Thricing replaces the kernel by a higher-order kernel. The same smoothing is the density and, after one integration, the call. A Gaussian damping of the characteristic function of the spline makes the convolution closed form. With the unquoted tails filled in, the estimator converges to the true density on the interior of a compact interval as the strike mesh fills that interval and the bandwidth tends to zero, and the same limit holds for noisy quotes when \(nh^{5}\to\infty\).

We compare the fit to several extant methods. \citeauthor{bondarenko2003}'s method is a sieve of fixed dimension: it does not become finer as more strikes arrive, and it is a density by construction. \citet{yatchew2006} replace the quoted calls by the nearest decreasing convex list; the second differences of that list are a discrete nonnegative measure, not a curve, so that projection is not in the comparison below. \citet{aitsahalia2003} smooth the same constrained list with a local polynomial, so the density is nonnegative by design and the call cannot match quotes that violate convexity. On the listed chains of Section~\ref{sec:emp} the out-of-the-money root-mean-square error (RMSE) is below the mixture of \citeauthor{bondarenko2003} and with the pricing-bandwidth kernels of \citet{aitsahalia1998} and \citet{grith2012}. The interpolant can lie below the \citeauthor{yatchew2006} floor because it is not required to be convex. 

The remainder of this manuscript is structured as follows. Section~\ref{sec:lit} reviews related work. Section~\ref{sec:est} constructs the estimator. The characteristic function of the normalized call supplies the closed form. Section~\ref{sec:refine} records two refinements: completing the tails and thricing. Section~\ref{sec:ifc} proves in-fill consistency for exact quotes and for noisy quotes. Section~\ref{sec:h} chooses bandwidths that meet those conditions and compares them on exact Heston and variance-gamma quotes. Section~\ref{sec:emp} presents an empirical illustration on four listed chains. Section~\ref{sec:conc} concludes.

\section{Related Work}
\label{sec:lit}

The identity of Section~\ref{sec:est} is the continuous-strike version of a spanning argument of \citet{ross1976}: if one could trade a put at every strike, the second derivative of the put curve would be a density, and one could synthesize any sufficiently regular payoff. \citet{dupire1994} obtained the analogous result in calendar time, which produces a local-volatility equation. Observed markets supply neither a continuum of strikes nor noiseless quotes.

A first family of methods specifies the density of \(S_T\) as a parametric curve and fits the implied option prices to the quotes. \citet{ritchey1990} and \citet{melick1997} take the density to be a mixture of lognormals, so that the call price is a mixture of Black--Scholes prices; \citet{bahra1997} and \citet{soderlind1997} use related specifications for interest-rate and foreign-exchange options. Two or three components typically suffice to generate a unimodal, skewed density, but the fit is poorly identified once more than two components are admitted, and the approximation to a Heston or jump-diffusion density need not vanish as more quotes arrive. Expansion methods perturb a baseline density---usually lognormal---by Gram--Charlier or Hermite terms \citep{jarrowrudd1982,corradosu1996,rompolis2007,xiu2014}. The expansions can match a finite number of moments, but they are not guaranteed to stay nonnegative, and the truncated series can oscillate in the tails.

A second, and in applied work the most common, family interpolates the Black--Scholes implied-volatility function of strike and converts the interpolant back to prices before differentiating twice. \citet{shimko1993} fitted a quadratic in strike; \citet{malz1997} used a cubic spline in delta; \citet{bliss2002} a smoothing spline. \citet{figlewski2010} completes the tails with a generalized extreme-value distribution. Because implied volatility is a much flatter function of strike than the call itself, a low-order interpolant often has small bias when quotes are accurate. The drawback is that implied volatility is a nonlinear transformation that inflates price noise where the derivative of price with respect to implied volatility is small, i.e.\ in the tails. An interpolating spline then treats those noisy tail quotes as exact, and the second derivative of the resulting call curve oscillates and can become negative. Shape-constrained smoothers of the implied-volatility surface \citep{fengler2005,fengler2009,gatheraljacquier2014,ludwig2015} address the sign problem at the cost of a more involved optimization.

A third family estimates the call, or its Black--Scholes implied volatility, by kernel regression or local polynomials and reads the density off a second strike derivative. \citet{aitsahalia1998} use a Nadaraya--Watson kernel. The state-price densities in that paper are estimated from \(14{,}431\) S\&P~500 options pooled across every trading day of 1993, because a single expiry was treated as too short a list on which to differentiate twice. They emphasise that the bandwidth dominates the choice of kernel: too small a neighborhood reproduces the noise in the quotes, and too large a neighborhood removes the curvature of the smile. \citet{aitsahalia2000} used the same estimator to recover an implied risk-aversion function. \citet{aitsahalia2003} work one day at a time, on the twenty to fifty strikes of a single expiry, and impose monotonicity and convexity of \(C\) in \(K\). The local polynomial that the asymptotics recommend for a second derivative is a local cubic \citep{fangijbels1996}. On \(242\) days in 1999, with the \(25\) most active strikes, unconstrained Nadaraya--Watson, local linear, local quadratic, and local cubic estimates of the density took a negative value on \(242\), \(130\), \(205\), and \(241\) days. The shape-constrained local linear never did. Stopping the violations in the unconstrained fits required substantial oversmoothing and a large bias. The raw quotes themselves broke convexity on \(24\%\) of the high-frequency prints that year.

\citet{grith2012} survey this literature and estimate a local cubic of the rescaled call with a quartic kernel. For a second derivative the optimal bandwidth is of order \(n^{-1/9}\). On their one-month DAX chain, least-squares cross-validation selects \(h=114\), which they describe as wiggly, so the densities they plot use \(h=228\) and \(h=434\). The implied-volatility smoother they display is widened further, to \(h=372\). Smoothing the call is more sensitive to that choice than smoothing implied volatility, and the tails oscillate unless the bandwidth is increased at a cost in bias. Twice differentiating a Nadaraya--Watson or \citet{gassermuller1979} level fit is set aside as the statistically worse route to the same derivative. \citet{jackwerth2004} starts a kernel regression of implied volatility from Silverman's rule, of order \(n^{-1/5}\). He warns that the bandwidth which prices overfits noise from one day to the next, that a kernel does not continue across gaps in the strike list, and that a spline has to be of degree higher than three if the recovered distribution is to be smooth. \citet{figlewski2018} finds that filling the implied-volatility curve with a cubic spline puts spikes in the density, because the second derivative of a cubic is continuous but not differentiable, and recommends a spline of degree four or higher. \citet{fanmancini2009} combine a parametric guide with a nonparametric correction. Related shape-constrained least squares are studied by \citet{yatchew2006} and \citet{birke2009}.

The estimator of Section~\ref{sec:est} is exposed to these objections. It is a kernel smooth of a single maturity. The pilot is a cubic spline of the complementary distribution function, the kernel is Gaussian, and the listed bandwidth is a fixed multiple of \(n^{-1/9}\), shorter than the constants used above to calm a local polynomial. Thricing raises the order of the kernel and puts negative sidelobes in the equivalent kernel for \(f_{\Q}\). The convolution is closed form. Without that correction, the smooth of the call is the fixed-design estimator of \citet{priestley1972} applied to a cubic interpolant, so that object is not entered as a separate method. Section~\ref{sec:emp} asks whether the listed fit still has to split the bandwidth, one width for the price and a wider one for a unimodal density. That section implements, at a pricing cross-validation bandwidth, the Nadaraya--Watson implied-volatility regression of \citeauthor{aitsahalia1998} and the local cubics of \citeauthor{grith2012}, together with the two-step procedure of \citeauthor{aitsahalia2003}: a least-squares projection of the quoted calls onto the decreasing convex cone, followed by a local polynomial of the constrained values whose pricing bandwidth is chosen the same way. The density drawn for \citeauthor{aitsahalia2003} uses the wider rule \(0.9\,F\hat\sigma\sqrt{T}\,n^{-1/9}\).

\citet{bondarenko2003} proposed the positive convolution approximation (PCA). The density is required to lie in the set of functions obtained by convolving a discrete probability with a positive kernel, and the discrete probability is chosen to minimize the sum of squared pricing errors. The kernel enforces smoothness and, together with positive weights, guarantees that the result is a density. In \citeauthor{bondarenko2003}'s comparison, the method dominated mixtures of lognormals, Hermite expansions, and several regularization methods. Subsequent simulation evidence on affine jump-diffusions is consistent with this ranking \citep{lu2019}. Section~\ref{sec:emp} implements a log-price analog: forty-one lognormals of common width, centered on a uniform grid in log strike from the first quoted strike to the last, fitted by nonnegative least squares. Sum-to-one and forward-mean penalties stay penalties; the weights are not rescaled after the fit. The width is chosen by even/odd pricing error.

\citet{rubinstein1994} and \citet{jackwerth1996} recovered probabilities on a binomial tree by minimizing a distance to a prior subject to pricing constraints. \citet{stutzer1996} and \citet{buchen1996} maximized entropy. These methods produce a valid density by construction, but the finite-support discretization and the choice of prior affect the tails. Viewing the problem as a linear inverse problem, \citet{cont2004} calibrated jump-diffusion models under a relative-entropy penalty; \citet{monnier2014} used the singular value decomposition of a restricted put operator.

A number of papers compare methods on market data, where the true density is unknown, so that the criterion is typically pricing error or stability over time \citep{cooper1999,jondeau2000,coutant2001,anagnou2005,bliss2002,liu2007}. Studies that generate options from a known process are less common. \citet{bondarenko2003} finds that his method outperforms mixtures, expansions, and several regularization methods on integrated squared error. \citet{lai2014} finds that kernel regression outperforms spline interpolation and neural nets on pricing error. \citet{santos2015} compare mixtures, Edgeworth expansions, hypergeometric functions, and a \citeauthor{bliss2002} smoothing spline against Heston and other targets. \citet{lu2019} compares methods on affine jump-diffusions and again finds \citeauthor{bondarenko2003}'s method accurate. The typical ranking is that implied-volatility interpolants price well when quotes are clean, that mixtures are unstable once more than two components are admitted, and that \citeauthor{bondarenko2003}'s method is the most accurate among methods that do not use a parametric density of \(S_T\), when the true density is known. Section~\ref{sec:heston} is in this vein: the targets are Heston and variance-gamma, and the true density is known. Section~\ref{sec:emp} scores listed index chains, where the criterion is pricing error, against \citeauthor{bondarenko2003}, \citeauthor{aitsahalia2003}, \citeauthor{aitsahalia1998}, and \citeauthor{grith2012}.

\section{Construction of the Estimator}
\label{sec:est}

\subsection{Construction}
\label{sec:construct}

Assume a risk-neutral measure \(\Q\) under which \(S_T\) has a density \(f_{\Q}\) on \((0,\infty)\) and \(\E^{\Q}[S_T]<\infty\). Let \(r\) be the continuously compounded interest rate and \(q\) the continuous dividend yield. Write \(F=S_0e^{(r-q)T}\) for the forward. Under \(\Q\),
\begin{align}
S_0e^{-qT}
&\;=\;
e^{-rT}\,\E^{\Q}[S_T]
\nonumber\\
&\;=\;
e^{-rT}\int_{0}^{\infty}x\,f_{\Q}(x)\,\mathrm{d}x,
\label{eq:S0}
\end{align}
hence \(\E^{\Q}[S_T]=F\). Write \(C\) for the corresponding call curve:
\begin{align}
C(K)
&\;=\;
e^{-rT}\,\E^{\Q}\bigl[(S_T-K)^+\bigr]
\nonumber\\
&\;=\;
e^{-rT}\int_K^{\infty}(x-K)\,f_{\Q}(x)\,\mathrm{d}x.
\label{eq:C}
\end{align}
Now define the normalized option price
\begin{equation}
\label{eq:Ctil}
\tilde C(K)
\;\equiv\;
\frac{C(K)}{S_0e^{-qT}}
\;=\;
\frac{\displaystyle\int_K^{\infty}(x-K)\,f_{\Q}(x)\,\mathrm{d}x}{F}.
\end{equation}

The normalized call price is a decreasing function of \(K\): it is equal to \(1\) if \(K=0\) and equal to \(0\) if \(K=\infty\). Thus the normalized call price, as a function of \(K\), is a complementary cdf, say \(\tilde C(K)=1-P(K)\). The characteristic function (hereafter c.f.) of the law of \(P\) is the Fourier--Stieltjes transform
\begin{equation}
\label{eq:phi-P-stieltjes}
\varphi_P(s)
\;=\;
\int_{(0,\infty)}e^{\mathrm{i}s K}\,\mathrm{d}P(K).
\end{equation}
For a cdf on \([0,\infty)\) and \(s\neq 0\), \eqref{eq:phi-P-stieltjes} is
\begin{equation}
\varphi_P(s)
\;=\;
1+\mathrm{i}s\int_0^{\infty}e^{\mathrm{i}s K}\,\tilde C(K)\,\mathrm{d}K.
\end{equation}
Substitute \eqref{eq:Ctil}. Fubini yields
\[
\int_0^{\infty}e^{\mathrm{i}s K}\,\tilde C(K)\,\mathrm{d}K
\;=\;
\frac{1}{F}\int_0^{\infty}f_{\Q}(x)\Biggl[\int_0^x e^{\mathrm{i}s K}(x-K)\,\mathrm{d}K\Biggr]\mathrm{d}x.
\]
The inner integral is \(-x/(\mathrm{i}s)+(e^{\mathrm{i}s x}-1)/(\mathrm{i}s)^2\). Hence, for \(s\neq 0\),
\begin{align}
\varphi_P(s)
&\;=\;
1+\frac{\mathrm{i}s}{F}\int_0^{\infty}f_{\Q}(x)\Biggl[-\frac{x}{\mathrm{i}s}+\frac{e^{\mathrm{i}s x}-1}{(\mathrm{i}s)^2}\Biggr]\mathrm{d}x
\nonumber\\
&\;=\;
1-F^{-1}\E^{\Q}[S_T]+\frac{\varphi_{\Q}(s)-1}{\mathrm{i}s\,F}
\nonumber\\
&\;=\;
\frac{\varphi_{\Q}(s)-1}{\mathrm{i}s\,F},
\label{eq:phi-P}
\end{align}
where \(\varphi_{\Q}(s)=\E^{\Q}[e^{\mathrm{i}s S_T}]\) and \(\E^{\Q}[S_T]=F\) by \eqref{eq:S0}. Solving for \(\varphi_{\Q}\), one obtains
\begin{equation}
\label{eq:phi-Q}
\varphi_{\Q}(s)
\;=\;
1+\mathrm{i}s\,F\,\varphi_P(s).
\end{equation}
Read from the right, \eqref{eq:phi-P} is a pricing formula: a model that supplies \(\varphi_{\Q}\) yields \(\varphi_P\), hence the complementary cdf \(\tilde C=1-P\) and the call \citep{lewis2000}. This paper inverts the identity: \(\varphi_P\) is constructed from the quotes, and \(\varphi_{\Q}\) is recovered.

The c.f.\ \(\varphi_{\Q}\) is that of \(S_T\). Fourier inversion yields
\begin{equation}
\label{eq:f-Q}
f_{\Q}(K)
\;=\;
\frac{1}{2\pi}\int_{-\infty}^{\infty}e^{-\mathrm{i}s K}\,\varphi_{\Q}(s)\,\mathrm{d}s.
\end{equation}
Substitute \eqref{eq:phi-Q} into \eqref{eq:f-Q}:
\begin{equation}
\label{eq:f-Q-plug}
f_{\Q}(K)
\;=\;
\frac{1}{2\pi}\int_{-\infty}^{\infty}e^{-\mathrm{i}s K}\,\mathrm{d}s
+\frac{F}{2\pi}\int_{-\infty}^{\infty}e^{-\mathrm{i}s K}\,(\mathrm{i}s)\,\varphi_P(s)\,\mathrm{d}s.
\end{equation}
The first integral is a Dirac mass at \(K=0\). For \(K>0\),
\begin{equation}
\label{eq:f-Q-pos}
f_{\Q}(K)
\;=\;
\frac{F}{2\pi}\int_{-\infty}^{\infty}e^{-\mathrm{i}s K}\,(\mathrm{i}s)\,\varphi_P(s)\,\mathrm{d}s.
\end{equation}

Observed calls at strikes \(0<K_1<\cdots<K_m\) supply values \(\tilde C_i=C_i/(S_0e^{-qT})\) and hence \(P_i=1-\tilde C_i\). The estimator interpolates these knots by a natural cubic spline \(P^{\mathrm{cs}}\). The spline matches value and first derivative at the knots. The c.f.\ of this absolutely continuous measure is \(\int e^{\mathrm{i}s K}\,P^{\mathrm{cs}\prime}(K)\,\mathrm{d}K\). Replace \(\varphi_P\) by the product of that integral with the Fourier transform of a Gaussian kernel of width \(h\),
\begin{equation}
\label{eq:phi-hat}
\hat\varphi_P(s)
\;=\;
\Biggl(\int_{(0,\infty)}e^{\mathrm{i}s K}\,P^{\mathrm{cs}\prime}(K)\,\mathrm{d}K\Biggr)e^{-h^2 s^2/2}.
\end{equation}
The Gaussian density and its distribution function are
\begin{equation}
\label{eq:kappa}
\kappa_h(u)
\;=\;
\frac{1}{h\sqrt{2\pi}}\,\exp\Bigl(-\frac{u^2}{2h^2}\Bigr),
\qquad
\mathcal{K}_h(u)
\;=\;
\int_{-\infty}^{u}\kappa_h.
\end{equation}
The factor \(e^{-h^2 s^2/2}\) is convolution against \(\kappa_h\), and multiplication by \(\mathrm{i}s\) is differentiation up to sign. Substituting \eqref{eq:phi-hat} into \eqref{eq:f-Q-pos} therefore yields
\begin{equation}
\label{eq:cubic}
\hat f_{\Q}(K)
\;=\;
-F\,(P^{\mathrm{cs}\prime}\ast\kappa_h)'(K),
\qquad
\hat P(K)
\;=\;
(P^{\mathrm{cs}\prime}\ast\mathcal{K}_h)(K).
\end{equation}
On each cell the integrand is a quadratic times a Gaussian. The integrals of \(1\), \(u\), and \(u^2\) against \(\phi\) and \(\Phi\) are elementary, so \eqref{eq:cubic} is closed form. Calls are \(\hat C=S_0e^{-qT}(1-\hat P)\). These knots are the quoted strikes. Section~\ref{sec:tails} adds the endpoints that complete the distribution.
\label{sec:closed}

\section{Two Refinements}
\label{sec:refine}

\subsection{Completing the Tails}
\label{sec:tails}

The spline of Section~\ref{sec:construct} stops at the first and last quotes. On \((0,K_1)\) it does not yet specify \(P\), and if \(C_m>0\) the residual mass \(1-P(K_m)\) is unplaced. Both ends are wrong as a cdf on \((0,\infty)\).

On \((0,K_1]\) a deep in-the-money call is intrinsic, \(C(K)=e^{-rT}(F-K)\), hence \(\tilde C(K)=1-K/F\) and \(P(K)=K/F\). That is a straight line through the origin. The estimator matches the observed \(P(K_1)\) rather than \(K_1/F\),
\begin{equation}
\label{eq:left}
P(K)
\;=\;
P(K_1)\,\frac{K}{K_1},
\qquad
0<K\le K_1.
\end{equation}
When the first quote is at intrinsic, \(P(K_1)=K_1/F\) and \eqref{eq:left} is exactly \(K/F\). The spline imposes the endpoint \(P(0)=0\), which is the same line at the boundary.

On the right, \(\tilde C(\infty)=0\) so \(P(\infty)=1\). If \(C_m>0\), a European call still has to fall to zero. The residual call is continued linearly with slope in \([-e^{-rT},0)\) until \(C=0\) at an intercept \(K_{\mathrm{end}}\). The slope is the last strictly negative step of the quoted call, clipped to that interval. A flat last step, two quotes on the premium floor, is not a decay: using it would send \(K_{\mathrm{end}}\) to the cap \(8F\). The spline imposes the single endpoint \(P(K_{\mathrm{end}})=1\). The bandwidth still counts a notional knot spacing along that line, eight per last quoted gap,
\begin{equation}
\label{eq:nR}
n_R
\;=\;
\mathrm{round}\bigl(8(K_{\mathrm{end}}-K_m)/(K_m-K_{m-1})\bigr),
\end{equation}
are equally spaced in \((K_m,K_{\mathrm{end}}]\) and converted to \(P=1-\tilde C\). If \(C_m=0\), \(P(K_m)=1\) already and this wing is empty.

Only the left completion changes the estimator on the Heston and variance-gamma chains below: the first quoted call is at intrinsic, and the last quoted call is already zero (dense) or of order \(10^{-5}\) (sparse).

\subsection{Thricing}
\label{sec:thrice}

The Gaussian kernel \(\kappa_h\) has second moment \(h^2\), so the leading bias of \eqref{eq:cubic} is \(O(h^2)\). When the quotes are treated as exact there is no observation-error variance to inflate, and a higher-order kernel is admissible. Twicing \citep{schucanysommers1977} replaces \(e^{-h^2 s^2/2}\) by \(2e^{-h^2 s^2/2}-e^{-h^2 s^2}\), the transform of \(2\kappa_h-\kappa_{h\sqrt{2}}\). That kernel has a vanishing second moment, and the bias falls from \(O(h^2)\) to \(O(h^4)\). One further bandwidth cancels the new leading term. The transform in \eqref{eq:phi-hat} is replaced by
\begin{equation}
\label{eq:phi-thrice}
\frac83 e^{-h^2 s^2/2}-2e^{-h^2 s^2}+\frac13 e^{-2h^2 s^2},
\end{equation}
which is the Fourier transform of \(\tfrac83\kappa_h-2\kappa_{h\sqrt{2}}+\tfrac13\kappa_{2h}\). The moments of order \(2\) and \(4\) vanish, so the bias is \(O(h^6)\) when \(f_{\Q}\) has a bounded sixth derivative. In closed form,
\begin{equation}
\label{eq:thrice}
\hat f_{\Q}^{\mathrm{th}}(K)
\;=\;
\frac83\hat f_{\Q,h}(K)-2\hat f_{\Q,h\sqrt{2}}(K)+\frac13\hat f_{\Q,2h}(K),
\end{equation}
and the same combination is applied to the interpolant in \eqref{eq:cubic}. The equivalent kernel takes negative values, so \(\hat f_{\Q}^{\mathrm{th}}\) is not a density. The call is \(S_0e^{-qT}(1-\hat P^{\mathrm{th}})\), with \(\hat P^{\mathrm{th}}\) clipped to \([0,1]\) and both premiums floored at zero.

Thricing a truncated distribution amplifies leakage from the missing tails, so it is applied only after Section~\ref{sec:tails}. The estimator used below is \eqref{eq:cubic} with those tails and with \eqref{eq:thrice}.

\section{In-Fill Consistency}
\label{sec:ifc}

Write \(P\) for the complementary-call cdf of Section~\ref{sec:est}, so that \(f_{\Q}=-F\,P''\) on \((0,\infty)\). The estimator \eqref{eq:cubic} replaces \(P'\) by the derivative of its cubic spline and smooths with \(\kappa_h\).

\begin{theorem}
\label{thm:ifc}
Let \(\E^{\Q}[S_T]<\infty\), and let \(f_{\Q}\) be four times continuously differentiable on \((0,\infty)\). Let \(C\) be the corresponding call curve. Fix a compact interval \([A,B]\subset(0,\infty)\). Consider a sequence of exact quote grids \(0<K_1^{(m)}<\cdots<K_m^{(m)}\) with mesh \(\delta_m=\max_i(K_i^{(m)}-K_{i-1}^{(m)})\to 0\), \(K_1^{(m)}\le A\) and \(K_m^{(m)}\ge B\). Form the cubic spline of Section~\ref{sec:construct} and complete the tails as in Section~\ref{sec:tails}. Let \(h_m\to 0\). Then both \eqref{eq:cubic} and its thriced version \eqref{eq:thrice} converge to \(f_{\Q}\) uniformly on every compact subset of \((A,B)\).
\end{theorem}

\begin{proof}
The estimator is \(\hat f_{\Q}=-F(P^{\mathrm{cs}\prime}\ast\kappa_h)'\), and \(f_{\Q}=-F P''\). Split
\[
(P^{\mathrm{cs}\prime}\ast\kappa_h)'-P''
\;=\;
\bigl((P^{\mathrm{cs}\prime}-P')\ast\kappa_h\bigr)'
+(P'\ast\kappa_h-P')'.
\]
The second term is ordinary kernel bias of \(P''=-f_{\Q}/F\). Because \(\kappa\) is even, \(\int\kappa=1\) and \(\int u\kappa(u)\,\mathrm{d}u=0\). If \(f_{\Q}''\) is bounded on a neighborhood of \([A,B]\), the difference is \(O(h^2)\) uniformly on compact subsets of \((A,B)\) several bandwidths away from the endpoints. The thriced kernel \(\tfrac83\kappa_h-2\kappa_{h\sqrt{2}}+\tfrac13\kappa_{2h}\) has Fourier transform \eqref{eq:phi-thrice}. Its moments of order \(2\) and \(4\) vanish, so the same expansion starts at order \(h^6\) whenever \(f_{\Q}\) is \(C^6\).

The first term is interpolation. On a quasi-uniform mesh, the natural cubic spline of a \(C^4\) function converges to the function, and its first derivative converges to the first derivative, uniformly on compact subsets of the interior of the grid. Thus \(P^{\mathrm{cs}\prime}\to P'\) as \(\delta_m\to 0\). Convolution against \(\kappa_h\) preserves the limit, and one derivative in \(K\) does not change the conclusion on those compacts.

The filled wings are supported on \((0,K_1]\) and, if \(C_m>0\), on \((K_m,K_{\mathrm{end}}]\). Their contribution at an evaluation point \(K\in[A+\varepsilon,B-\varepsilon]\) is at most a constant times \(h^{-2}\exp(-\varepsilon^2/2h^2)\), because \(K_1\le A\) and \(K_m\ge B\). That bound tends to zero faster than any power of \(h\).

Multiplying by \(-F\) yields the claim for \(\hat f_{\Q}\) and for \(\hat f_{\Q}^{\mathrm{th}}\).
\end{proof}

\begin{theorem}
\label{thm:noise}
Keep the assumptions of Theorem~\ref{thm:ifc} on \(f_{\Q}\) and on the grids, and assume in addition that the mesh on \([A,B]\) is quasi-uniform: \(n_m\delta_m\) is bounded above and below, where \(n_m\) is the number of strikes in \([A,B]\). Observe
\[
P_i^{(m)}
\;=\;
P(K_i^{(m)})+\eta_i^{(m)},
\]
with \(\eta_i^{(m)}\) independent, \(\E[\eta_i^{(m)}]=0\) and \(\E[(\eta_i^{(m)})^2]\le\tau^2<\infty\). Form the cubic spline and the tails from these noisy knots. Let \(h_m\to 0\) and \(n_m h_m^5\to\infty\). Then both \eqref{eq:cubic} and \eqref{eq:thrice} converge in probability to \(f_{\Q}\), uniformly on every compact subset of \((A,B)\).
\end{theorem}

\begin{proof}
The spline operator is linear, so the spline of the noisy knots is the spline of the true knots plus the spline \(S_\eta\) of \((K_i^{(m)},\eta_i^{(m)})\), up to the tail knots. The tail knots lie outside \([A,B]\). Their contribution on a compact several bandwidths inside \((A,B)\) is of order \(h^{-2}\exp(-c/h^2)\) as in the proof of Theorem~\ref{thm:ifc}, whether or not the intercept \(K_{\mathrm{end}}\) depends on the last noisy slope, because that intercept stays inside a fixed multiple of \(F\).

The piece built from the true knots converges to \(f_{\Q}\) by Theorem~\ref{thm:ifc}. The noise piece is \(-F(S_\eta'\ast\kappa_h)'\). On a quasi-uniform mesh, once \(h/\delta\to\infty\), this is within \(o_p(1)\) of the kernel estimator
\[
\sum_{i=1}^{n_m}\eta_i^{(m)}\,\delta_i\,h_m^{-3}L\bigl((K-K_i^{(m)})/h_m\bigr),
\]
where \(\delta_i\) is the local gap and \(L\) is a fixed linear combination of \(\kappa''\) (of \(\tfrac83\kappa''(\cdot)-2\kappa''(\cdot/\sqrt{2})+\tfrac13\kappa''(\cdot/2)\) after thricing). The kernel \(L\) is bounded and integrable, and \(\int L=0\). The sum has conditional mean zero. Its conditional variance is of order
\[
\tau^2\sum_i\delta_i^2 h_m^{-6}L^2\bigl((K-K_i)/h_m\bigr)
\;\asymp\;
\tau^2\delta_m h_m^{-5}\int L^2
\;=\;
O\bigl(\tau^2/(n_m h_m^5)\bigr),
\]
uniformly for \(K\) in a compact subset of \((A,B)\). The hypothesis \(n_m h_m^5\to\infty\) sends this variance to zero, and it also forces \(h_m/\delta_m=n_m h_m\to\infty\). Chebyshev's inequality gives pointwise convergence in probability. The map \(K\mapsto(S_\eta'\ast\kappa_h)'(K)\) is Lipschitz with a random constant of that same stochastic order, so a grid of spacing \(h_m\) upgrades the convergence to uniform convergence on the compact.

Thricing replaces \(L\) by a fixed linear combination of three such kernels and does not change the order of the variance. The bias expansion remains the one in the proof of Theorem~\ref{thm:ifc}.
\end{proof}

\begin{remark}
Theorem~\ref{thm:ifc} is a statement about exact quotes. Theorem~\ref{thm:noise} is the corresponding statement about additive noise in \(P\). Neither is a statement about the far tails outside \([A,B]\). The equivalent kernel of \eqref{eq:thrice} takes negative values, so the thriced estimator is not necessarily a density. The empirical section fits the spline to the decreasing-convex projection of the quoted calls. When the true call lies in the interior of that cone, the projection changes the knots by an amount of the same stochastic order as \(\eta_i\), and the bandwidth condition of Theorem~\ref{thm:noise} still applies.
\end{remark}

\section{Choice of Bandwidth}
\label{sec:h}

The width \(h\) is the only tuning constant in \eqref{eq:thrice}. Two rules are used below. Each is checked against Theorems~\ref{thm:ifc} and \ref{thm:noise}.

With exact quotes the bias is the usual \(O(h^2)\) kernel bias, or \(O(h^6)\) after thricing, and the spline error vanishes with the mesh. For a Gaussian it is enough that \(h\) be a small multiple of the typical quoted gap. Write \(\delta\) for the median of \(\{K_i-K_{i-1}\}\) on the quoted strikes, not on the filled wings, and take
\begin{equation}
\label{eq:h-mesh}
h
\;=\;
\min\bigl(1.2\,\delta,\; 0.30\,(K_m-K_1)\bigr).
\end{equation}
The cap binds only on a very short chain. On an equally spaced grid every gap equals \(\delta\). Whenever \(\delta\to 0\), this rule satisfies \(h\to 0\), which is what Theorem~\ref{thm:ifc} requires. It does not satisfy \(nh^5\to\infty\). A noisy chain with mesh \(\delta\) and \(h\propto\delta\) has \(nh^5\) of order \(\delta^4\to 0\), so the variance in the proof of Theorem~\ref{thm:noise} explodes. The exponent \(n^{-1/5}\) familiar from kernel density estimation leaves the same product bounded, so it is not used.

The density is a second derivative. The bandwidth that balances an \(O(h^2)\) bias against a variance of order \(1/(nh^5)\) is of order \(n^{-1/9}\), the rate \citet{aitsahalia2003} use for their local cubic. Any fixed positive multiple satisfies both theorems: \(h\to 0\) and \(nh^5\to\infty\). Write \(\hat\sigma_{\mathrm{ATM}}\) for the at-the-money Black--Scholes implied volatility, the median of contracts with \(|K-F|\le 0.03F\) (or the median of all finite implied volatilities if that window is empty), and take
\begin{equation}
\label{eq:h-d2}
h
\;=\;
c\,F\,\hat\sigma_{\mathrm{ATM}}\sqrt{T}\,n^{-1/9},
\end{equation}
with \(n\) the number of quoted strikes and \(c=0.35\). On chains of one to several hundred strikes this is the bandwidth used for the listed prices of our estimator. Section~\ref{sec:heston} reports both rules.

The exponent is the one \citet{grith2012} derive for a local cubic of a second derivative and the order \citet{aitsahalia2003} take from \citet{fangijbels1996}. The bandwidths those papers plot are wider. \citeauthor{grith2012} double and quadruple a cross-validated width because the cross-validated density is wiggly, and \citeauthor{aitsahalia2003} report that an unconstrained local polynomial stops taking negative values only after substantial oversmoothing. \citet{jackwerth2004} starts from a multiple of \(n^{-1/5}\). That exponent leaves \(nh^{5}\) bounded, which is why it is not used. Any fixed \(c>0\) sends \(nh^{5}\) to infinity. The constant \(c=0.35\) is the short end of that family. On the four listed chains it is just under two-fifths of the factor \(0.9\) used for the \citeauthor{aitsahalia2003} density in Section~\ref{sec:emp}.

\subsection{Numerical Illustration}
\label{sec:heston}

The data-generating process is the stochastic-volatility model of \citet{heston1993} with the equity parameters of \citet{bakshi1997}: \(S_0=100\), \(r=0.05\), \(q=0\), \(T=1/4\), \(\kappa=1.15\), \(\theta=0.04\), \(\sigma=0.39\), \(\rho=-0.64\), \(v_0=0.04\). The c.f.\ of \(\log S_T\) uses the little-trap branch of \citet{albrecher2007}. Call prices are the damped integrand of \citet{carrmadan1999} with damping \(\alpha=1.5\), evaluated by the trapezoidal rule on \(4096\) frequencies in \((0,250]\); puts follow by parity. The true density is Fourier inversion of the same c.f.\ on the same frequency grid,
\begin{equation}
\label{eq:true-q}
f_{\Q}(K)
\;=\;
\frac1{K\pi}\int_0^\infty
\mathrm{Re}\bigl[e^{-\mathrm{i}u\log K}\varphi(u)\bigr]\,\mathrm{d}u,
\end{equation}
where \(\varphi(u)=\E^{\Q}[e^{\mathrm{i}u\log S_T}]\). Two chains are used: a dense grid of \(m=256\) equally spaced strikes on \([30,220]\) (spacing \(\delta=0.745\)) and a sparse grid of \(m=32\) equally spaced strikes on \([70,140]\) (spacing \(\delta=2.26\)). Integrated squared error (ISE) is \(\int_{60}^{150}(\hat f_{\Q}-f_{\Q})^2\,\mathrm{d}K\), evaluated by the trapezoidal rule on \(401\) equally spaced points. Table~\ref{tab:heston-tail} reports ISE for the three nested estimators at the oracle \(h^{\ast}\), at the mesh rule \eqref{eq:h-mesh}, and at \eqref{eq:h-d2}. The oracle minimizes the integral on a geometric grid in \(h\).

\begin{table}[H]
\centering
\caption{Integrated squared error on exact Heston quotes for the cubic-spline estimator. Mesh \(h\) is \eqref{eq:h-mesh}. The last pair is \eqref{eq:h-d2}. Quoted spline uses only the observed strikes. Tails adds the wings of Section~\ref{sec:tails}. Tails+Thrice applies \eqref{eq:thrice}.}
\label{tab:heston-tail}
\resizebox{\textwidth}{!}{%
\begin{tabular}{ll
  S[table-format=2.2]
  S[table-format=1.2e-1, scientific-notation=true]
  S[table-format=2.2]
  S[table-format=1.2e-1, scientific-notation=true]
  S[table-format=2.2]
  S[table-format=1.2e-1, scientific-notation=true]}
\toprule
\multicolumn{1}{c}{Chain} & \multicolumn{1}{c}{Estimator} &
  {\(h^{\ast}\)} & {Oracle ISE} & {Mesh \(h\)} & {Mesh ISE} & {\(h_{1/9}\)} & {\(h_{1/9}\) ISE} \\
\midrule
Dense & Quoted spline
  & 0.15 & 9.76e-10 & 0.89 & 7.36e-7 & 1.84 & 1.25e-5 \\
Dense & Tails
  & 0.15 & 9.76e-10 & 0.89 & 7.36e-7 & 1.84 & 1.25e-5 \\
Dense & Tails+Thrice
  & 0.63 & 1.16e-12 & 0.89 & 7.12e-12 & 1.84 & 1.16e-8 \\
\midrule
Sparse & Quoted spline
  & 16.19 & 1.80e-2 & 2.71 & 1.04e-1 & 2.31 & 1.22e-1 \\
Sparse & Tails
  & 0.45 & 1.13e-6 & 2.71 & 5.29e-5 & 2.31 & 2.94e-5 \\
Sparse & Tails+Thrice
  & 2.18 & 9.19e-7 & 2.71 & 1.14e-6 & 2.31 & 9.35e-7 \\
\bottomrule
\end{tabular}%
}
\end{table}

\begin{figure}[H]
\centering
\includegraphics[width=0.98\textwidth]{figures/ccdf_heston_rnd.pdf}
\caption{Heston density of \(S_T\) on \([60,150]\). The three estimators of Table~\ref{tab:heston-tail}, each at the mesh bandwidth \eqref{eq:h-mesh}. Sparse quoted strikes are marked on the true density. Curves are truncated at zero, and the vertical scale is that of the true density.}
\label{fig:heston-rnd}
\end{figure}

On the dense grid the first strike lies left of the window and the last call is already zero, so the wings do not move the error. The spline without thricing has mesh integrated squared error \(7.36\times 10^{-7}\). Thricing cuts that to \(7.12\times 10^{-12}\), against an oracle \(1.16\times 10^{-12}\). Equation \eqref{eq:h-d2} is wider (\(h=1.84\)) and the thriced error is \(1.16\times 10^{-8}\). On the sparse grid the quoted spline without tails is useless: mesh integrated squared error \(0.104\). The left line \eqref{eq:left} brings the plain mesh error to \(5.29\times 10^{-5}\). Thricing then reaches \(1.14\times 10^{-6}\) at the mesh and \(9.35\times 10^{-7}\) at \eqref{eq:h-d2} (\(h=2.31\)), against an oracle \(9.19\times 10^{-7}\). Figure~\ref{fig:heston-rnd} plots the three estimators at the mesh bandwidth.

Table~\ref{tab:vg-tail} repeats the design for the variance-gamma process of \citet{madan1998}. The process parameters are those of \citet{carrmadan1999}: \(\sigma=0.12\), \(\nu=0.2\), \(\theta=-0.14\). Spot, rate, dividend yield, and maturity are as above, as are the two strike grids and the ISE window. Variance gamma is a pure-jump L\'evy process. The true density of \(\log S_T\) is a modified Bessel function \citep{madan1998}; the density of \(S_T\) is that formula times \(1/K\). Calls are the discounted payoff of that density, evaluated by the trapezoidal rule.

\begin{table}[H]
\centering
\caption{Integrated squared error on exact variance-gamma quotes for the cubic-spline estimator. Columns as in Table~\ref{tab:heston-tail}.}
\label{tab:vg-tail}
\resizebox{\textwidth}{!}{%
\begin{tabular}{ll
  S[table-format=2.2]
  S[table-format=1.2e-1, scientific-notation=true]
  S[table-format=2.2]
  S[table-format=1.2e-1, scientific-notation=true]
  S[table-format=2.2]
  S[table-format=1.2e-1, scientific-notation=true]}
\toprule
\multicolumn{1}{c}{Chain} & \multicolumn{1}{c}{Estimator} &
  {\(h^{\ast}\)} & {Oracle ISE} & {Mesh \(h\)} & {Mesh ISE} & {\(h_{1/9}\)} & {\(h_{1/9}\) ISE} \\
\midrule
Dense & Quoted spline
  & 0.17 & 5.12e-7 & 0.89 & 5.08e-5 & 1.18 & 1.20e-4 \\
Dense & Tails
  & 0.17 & 5.12e-7 & 0.89 & 5.08e-5 & 1.18 & 1.20e-4 \\
Dense & Tails+Thrice
  & 0.46 & 4.04e-7 & 0.89 & 4.90e-6 & 1.18 & 1.43e-5 \\
\midrule
Sparse & Quoted spline
  & 14.03 & 3.27e-2 & 2.71 & 1.05e-1 & 1.46 & 1.93e-1 \\
Sparse & Tails
  & 0.45 & 3.97e-5 & 2.71 & 1.25e-3 & 1.46 & 2.43e-4 \\
Sparse & Tails+Thrice
  & 0.92 & 3.13e-5 & 2.71 & 3.00e-4 & 1.46 & 4.67e-5 \\
\bottomrule
\end{tabular}%
}
\end{table}

\begin{figure}[H]
\centering
\includegraphics[width=0.98\textwidth]{figures/ccdf_vg_rnd.pdf}
\caption{Variance-gamma density of \(S_T\) on \([60,150]\). The three estimators of Table~\ref{tab:vg-tail}, each at the mesh bandwidth \eqref{eq:h-mesh}. Sparse quoted strikes are marked on the true density. Curves are truncated at zero, and the vertical scale is that of the true density.}
\label{fig:vg-rnd}
\end{figure}

On the dense chain the wings again do not move the digits. Thricing cuts the mesh error from \(5.08\times 10^{-5}\) to \(4.90\times 10^{-6}\), against an oracle \(4.04\times 10^{-7}\). At \eqref{eq:h-d2} the thriced error is \(1.43\times 10^{-5}\) (\(h=1.18\)). The density of \(S_T\) is more peaked than under Heston, and the absolute errors are larger. On the sparse chain the quoted spline without tails has mesh error \(0.105\). Tails bring the plain mesh error to \(1.25\times 10^{-3}\), and thricing lowers it to \(3.00\times 10^{-4}\) at the mesh and \(4.67\times 10^{-5}\) at \eqref{eq:h-d2} (\(h=1.46\)), against an oracle \(3.13\times 10^{-5}\). Here the mesh bandwidth \(h=2.71\) is the wider of the two. Figure~\ref{fig:vg-rnd} plots the three estimators at the mesh bandwidth.

\section{Empirical Applications}
\label{sec:emp}

\subsection{Data and Methodology}
\label{sec:data}

The listed comparison uses European, AM-settled index options from Cboe's delayed-quotes feed (\texttt{cdn.cboe.com}). Each file is dated 23~September 2026. The S\&P~500 (SPX) file is 03:56 Eastern (30{,}470 listed contracts across 58 expiries); the Nasdaq-100 (NDX) file is 03:44 Eastern (16{,}492 contracts across 46 expiries); the Russell~2000 (RUT) file is 03:44 Eastern (10{,}790 contracts across 30 expiries). Only the standard monthly roots SPX, NDX, and RUT are kept, not weekly SPXW, NDXW, or RUTW. Time to expiry is \(T=d/365.25\) with \(d\) the calendar day count from the snapshot date. The discount rate is \(r=0.04\).

A quote is retained if the bid is positive, the ask is at least the bid, the mid is at least \(0.25\), and the spread is at most \(\max(15,\,0.002 S_0,\,0.5\times\mathrm{mid})\). The \(15\)-point term is the tolerance on the S\&P~500. On the Nasdaq-100 a fixed \(15\)-point cap deletes puts whose spreads are a dollar or two wider and leaves a gap of several thousand points in the strike list. Among remaining put--call pairs, the forward is the median of \(K+e^{rT}(C-P)\) over strikes with \(|K-K_{\mathrm{atm}}|\le 50\) on SPX, where \(K_{\mathrm{atm}}\) is the pair closest to the cash index; on NDX and RUT the band is \(0.65\%\) of the spot. Dividend yield is \(q=r-\log(F/S_0)/T\). The working slice is out of the money (OTM): the put mid if \(K\le F\), the call mid if \(K\ge F\), the missing side filled by parity. Puts with \(K<0.7F\) and mid below \(0.50\) are then dropped; out-of-the-money calls with mid at least \(0.25\) are kept. Table~\ref{tab:data} records the filter counts. In \eqref{eq:cubic} and \eqref{eq:thrice}, \(\tilde C=C/(S_0e^{-qT})\) and \(F=S_0e^{(r-q)T}\).

The December 2026 S\&P~500 AM quarterly (\(T=0.235\), \(S_0=7764.64\), \(F=7836\), \(q=0.11\%\)) appears in Figure~\ref{fig:listed}. After the screens it has \(n=444\) strikes on \([0.36F,1.30F]\), of which \(75\%\) are puts. Near the money the strike mesh is \(5\) (\(0.06\%\) of \(F\)); it is \(5\)--\(25\) on the put wing and up to \(200\) beyond \(1.2F\). Median out-of-the-money mids are \(39\) for puts and \(62\) for calls. Implied volatility \(\hat\sigma_{\mathrm{ATM}}=0.134\) is the median Black--Scholes implied volatility of contracts with \(|K-F|\le 0.03F\). The linear call tail of Section~\ref{sec:tails} is appended (\(C_m=0.25\) by the mid floor of the filter, \(n_R=80\) from \eqref{eq:nR}). The bandwidth \eqref{eq:h-d2} uses the \(444\) quoted strikes and equals \(h=90.9\). The March 2027 AM quarterly (\(T=0.485\), \(F=7918\), \(q=-0.04\%\), \(\hat\sigma_{\mathrm{ATM}}=0.144\)) filters to \(n=221\) strikes on \([0.20F,1.49F]\), with a \(25\)-point mesh at the money; \eqref{eq:nR} gives \(n_R=41\), and \eqref{eq:h-d2} gives \(h=152.8\).

The Nasdaq-100 December 2026 AM quarterly (\(T=0.235\), \(S_0=30732\), \(F=31025\), \(q=-0.03\%\), \(\hat\sigma_{\mathrm{ATM}}=0.201\)) is built with the same screens. It filters to \(n=312\) strikes on \([0.29F,1.39F]\), \(83\%\) puts; the at-the-money mesh is \(25\), and gaps of \(1000\) appear beyond \(1.2F\). The bandwidth \eqref{eq:h-d2} is \(h=558.2\) (\(n=312\); the right wing has \(n_R=72\)).

The Russell~2000 December 2026 AM quarterly (\(T=0.235\), \(S_0=2889.92\), \(F=2912\), \(q=0.80\%\), \(\hat\sigma_{\mathrm{ATM}}=0.183\)) filters to \(n=99\) strikes on \([0.45F,1.32F]\), \(59\%\) puts; the at-the-money mesh is \(10\). The last two calls are both \(0.25\), so the wing slope is the previous negative step. Equation \eqref{eq:nR} gives \(n_R=40\). The bandwidth \eqref{eq:h-d2} is \(h=54.4\).

The quoted calls are replaced by the nearest vector in Euclidean distance that is decreasing and convex in strike. Consecutive slopes lie in \([-e^{-rT},0]\), and each coordinate is at least the intrinsic value \(e^{-rT}\max(F-K_i,0)\).  This is the unpenalized projection of \citeauthor{yatchew2006} and also the first step of \citeauthor{aitsahalia2003}. Puts of the projected call follow by parity, and both premiums are floored at zero.

Out-of-the-money RMSE uses the quoted mid of the put for \(K\le F\) and that of the call for \(K\ge F\). Hold-out is the average of two even/odd splits (train on even strikes and score odd, then the reverse), with \(h\) recomputed from the training \(n\) and the projection recomputed from the training strikes only.
\begin{table}[H]
\centering
\caption{Construction of the working slices. ``At expiry'' counts listed puts and calls on that root and maturity with a positive bid and ask at least the bid. ``Screens'' applies the mid and spread filters. ``Working OTM'' is the slice used by the estimators.}
\label{tab:data}
{\small
\begin{tabular}{@{}l *{5}{S[table-format=3.0]}@{}}
\toprule
\multicolumn{1}{@{}c}{Slice} & {At expiry} & {Screens} & {Unique \(K\)} & {Working OTM} & {Puts} \\
\midrule
SPX 18 Dec 2026 & 927 & 922 & 473 & 444 & 335 \\
SPX 19 Mar 2027 & 453 & 451 & 228 & 221 & 165 \\
NDX 18 Dec 2026 & 631 & 631 & 319 & 312 & 259 \\
RUT 18 Dec 2026 & 218 & 218 & 115 & 99 & 58 \\
\bottomrule
\end{tabular}
}
\end{table}

\subsection{Estimators}
\label{sec:emp-est}

\paragraph{Ours}
Start from the projected calls of Section~\ref{sec:data}. Set \(P_i=1-C_i/(S_0e^{-qT})\). Pass a natural cubic spline through those values, with the tail knots of Section~\ref{sec:tails}, and invert at \eqref{eq:h-d2} with thricing. The listed prices are that interpolant, floored at zero, with puts by parity.

\paragraph{\citeauthor{aitsahalia2003}}
Start from the projected calls of Section~\ref{sec:data}. The call used for pricing is a local linear regression of those values, with Gaussian weights. Its bandwidth is chosen by even/odd pricing error from the grid \(0.25,0.5,1,2,4\) times \(F\hat\sigma\sqrt{T}\,n^{-1/5}\). The density drawn in Figure~\ref{fig:listed} is \(e^{rT}\) times twice the quadratic coefficient of a local cubic of the same list, at the fixed second-derivative bandwidth \(0.9\,F\hat\sigma\sqrt{T}\,n^{-1/9}\). Prices in the table are the local-linear calls, floored, with puts by parity. 

\paragraph{\citeauthor{bondarenko2003} PCA}
Start from the projected calls of Section~\ref{sec:data}. The density is a mixture of \(41\) lognormals with a common log-width \(h\) and nonnegative weights. The centers are equally spaced in log strike from the first quoted strike to the last. Each component prices a call in closed form, Black--Scholes with forward \(e^{z_j+h^2/2}\). The weights minimise the sum of squared errors in those projected calls, plus two penalties, one for \(\sum_j a_j=1\) and one for \(\sum_j a_j e^{z_j+h^2/2}=F\). Nonnegative least squares returns the weights. They are not divided by their sum afterwards; a weight above one is clipped to one. The width \(h\) is chosen by even/odd pricing error on a grid of multiples of the center spacing. Prices in Table~\ref{tab:listed} are those closed-form calls, floored, with puts by parity.

\paragraph{\citeauthor{aitsahalia1998}}
Start from the projected calls of Section~\ref{sec:data}. Regress Black--Scholes implied volatility on strike by Nadaraya--Watson with a Gaussian kernel. The bandwidth minimizes the same even/odd pricing error. Prices are the Black--Scholes calls at the smoothed volatility, floored, with puts by parity. The density is the strike derivative of that call, including the slope and curvature of the smoothed smile. The selected bandwidths are \(h=20\), \(50\), \(200\), and \(20\). The resulting curves have \(32\), \(14\), \(26\), and \(15\) peaks.

\paragraph{\citeauthor{grith2012}}
Start from the projected calls of Section~\ref{sec:data}. Both rows are a local cubic with a quartic kernel, and the bandwidth minimizes the same even/odd pricing error. The call row reads the density as \(e^{rT}\) times twice the quadratic coefficient of a local cubic of the call. Its bandwidths are \(h=37.8\), \(67.6\), \(239\), and \(40\), and the curves have \(3\), \(3\), \(13\), and \(1\) peaks. The implied-volatility row is the local polynomial of the smile used by \citeauthor{grith2012}: a local cubic of implied volatility, Black--Scholes prices at the local level, and a density that uses the local slope and curvature. Its bandwidths are \(h=20\), \(50\), \(200\), and \(52.8\), and the curves have \(22\), \(6\), \(12\), and \(2\) peaks. The call cubic is drawn on the left of Figure~\ref{fig:listed}. Figure~\ref{fig:listed-kern} draws the implied-volatility cubic together with \citeauthor{aitsahalia1998}. Prices in both rows are floored, with puts by parity.

Hold-out refits the projection and every tuned method on the training strikes only, including the choice of its tuning constant, and scores the other half against the quoted mids. The two directions are averaged. 

\begin{table}[p]
\centering
\caption{Out-of-the-money repricing. Every fit starts from the projection in Section~\ref{sec:data} and is scored on the quoted mids. Ours is the cubic-spline inversion at \eqref{eq:h-d2} (\(h=90.9\), \(152.8\), \(558.2\), \(54.4\)). \citeauthor{aitsahalia2003} is priced from the local linear at its cross-validated bandwidth; the mass is the local cubic at \(0.9\,F\hat\sigma\sqrt{T}\,n^{-1/9}\). \citeauthor{aitsahalia1998} and both \citeauthor{grith2012} rows use even/odd pricing bandwidths, and the mass is the density of that same fit. PCA is priced from the lognormal mixture. OTM is put RMSE for \(K\le F\) and call RMSE for \(K\ge F\), premiums floored at zero. Mass is \(\int\max(\hat q,0)\) on \([\max(50,0.2F),2.4F]\). TV is the total variation of \(\max(\hat q,0)\) on \([0.55F,1.40F]\), divided by twice the maximum; a unimodal curve scores \(1.0\). Hold-out refits the projection and the tuning on each half. Average is the equal-weight mean of the four chains. Bold method names are the unimodal curves in Figure~\ref{fig:listed}: the spline, the lognormal mixture, and \citeauthor{aitsahalia2003}.}
\label{tab:listed}
{\small
\setlength{\tabcolsep}{3pt}%
\renewcommand{\arraystretch}{0.9}%
\resizebox{\textwidth}{!}{%
\begin{tabular}{ll *{4}{S[table-format=3.3]} S[table-format=2.1] S[table-format=2.1]}
\toprule
\multicolumn{1}{c}{Chain} & \multicolumn{1}{c}{Method} &
  \multicolumn{1}{c}{RMSE OTM} & \multicolumn{1}{c}{Puts} &
  \multicolumn{1}{c}{Calls} & \multicolumn{1}{c}{Hold-out} &
  \multicolumn{1}{c}{Mass} & \multicolumn{1}{c}{TV} \\
\midrule
SPX Dec & \textbf{Ours}
  & 0.035 & 0.033 & 0.039 & 0.039 & 1.0 & 1.0 \\
SPX Dec & \textbf{PCA}
  & 0.038 & 0.035 & 0.046 & 0.040 & 1.0 & 1.0 \\
SPX Dec & \textbf{\citeauthor{aitsahalia2003}}
  & 0.337 & 0.252 & 0.518 & 0.444 & 1.0 & 1.0 \\
SPX Dec & \citeauthor{grith2012}, IV
  & 0.030 & 0.029 & 0.032 & 0.037 & 4.3 & 3.2 \\
SPX Dec & \citeauthor{grith2012}, call
  & 0.031 & 0.030 & 0.036 & 0.039 & 1.0 & 1.2 \\
SPX Dec & \citeauthor{aitsahalia1998}
  & 0.064 & 0.061 & 0.073 & 0.113 & 1.6 & 10.9 \\
\midrule
SPX Mar & \textbf{Ours}
  & 0.050 & 0.047 & 0.059 & 0.058 & 1.0 & 1.0 \\
SPX Mar & \textbf{PCA}
  & 0.079 & 0.063 & 0.115 & 0.079 & 1.0 & 1.0 \\
SPX Mar & \textbf{\citeauthor{aitsahalia2003}}
  & 0.596 & 0.401 & 0.963 & 0.789 & 1.0 & 1.0 \\
SPX Mar & \citeauthor{grith2012}, IV
  & 0.027 & 0.024 & 0.034 & 0.042 & 1.1 & 1.4 \\
SPX Mar & \citeauthor{grith2012}, call
  & 0.031 & 0.029 & 0.038 & 0.061 & 1.0 & 1.1 \\
SPX Mar & \citeauthor{aitsahalia1998}
  & 0.073 & 0.064 & 0.094 & 0.158 & 1.2 & 5.0 \\
\midrule
NDX Dec & \textbf{Ours}
  & 1.042 & 1.088 & 0.780 & 1.099 & 1.0 & 1.0 \\
NDX Dec & \textbf{PCA}
  & 1.147 & 1.131 & 1.221 & 1.187 & 1.0 & 1.0 \\
NDX Dec & \textbf{\citeauthor{aitsahalia2003}}
  & 2.499 & 2.418 & 2.858 & 3.278 & 1.0 & 1.0 \\
NDX Dec & \citeauthor{grith2012}, IV
  & 0.874 & 0.927 & 0.548 & 1.018 & 1.0 & 2.7 \\
NDX Dec & \citeauthor{grith2012}, call
  & 0.892 & 0.943 & 0.579 & 1.087 & 1.0 & 2.1 \\
NDX Dec & \citeauthor{aitsahalia1998}
  & 1.734 & 1.648 & 2.104 & 2.623 & 1.3 & 11.0 \\
\midrule
RUT Dec & \textbf{Ours}
  & 0.025 & 0.025 & 0.024 & 0.030 & 1.0 & 1.0 \\
RUT Dec & \textbf{PCA}
  & 0.034 & 0.030 & 0.039 & 0.038 & 1.0 & 1.0 \\
RUT Dec & \textbf{\citeauthor{aitsahalia2003}}
  & 0.339 & 0.293 & 0.394 & 0.467 & 1.1 & 1.0 \\
RUT Dec & \citeauthor{grith2012}, IV
  & 0.023 & 0.023 & 0.024 & 0.031 & 3.5 & 1.7 \\
RUT Dec & \citeauthor{grith2012}, call
  & 0.022 & 0.022 & 0.022 & 0.061 & 1.1 & 1.2 \\
RUT Dec & \citeauthor{aitsahalia1998}
  & 0.098 & 0.124 & 0.035 & 0.201 & 1.6 & 7.6 \\
\midrule
Average & \textbf{Ours}
  & 0.288 & 0.298 & 0.225 & 0.306 & 1.0 & 1.0 \\
Average & \textbf{PCA}
  & 0.325 & 0.315 & 0.355 & 0.336 & 1.0 & 1.0 \\
Average & \textbf{\citeauthor{aitsahalia2003}}
  & 0.943 & 0.841 & 1.183 & 1.245 & 1.0 & 1.0 \\
Average & \citeauthor{grith2012}, IV
  & 0.238 & 0.251 & 0.159 & 0.282 & 2.5 & 2.2 \\
Average & \citeauthor{grith2012}, call
  & 0.244 & 0.256 & 0.169 & 0.312 & 1.0 & 1.4 \\
Average & \citeauthor{aitsahalia1998}
  & 0.492 & 0.475 & 0.577 & 0.774 & 1.4 & 8.6 \\
\bottomrule
\end{tabular}}%
}
\end{table}

On SPX December the quoted calls sit \(0.022\) from the decreasing-convex cone, so a fit that stays inside the cone cannot undercut that error. The cross-validated local linear of \citeauthor{aitsahalia2003} is \(0.337\), \(0.596\), \(2.499\), and \(0.339\). With centers on the quoted log-strike range and a width chosen by pricing error, PCA out-of-the-money error is \(0.038\) on SPX December, \(0.079\) on SPX March, \(1.147\) on NDX December, and \(0.034\) on RUT December.

The estimator of this paper is the Gaussian inversion of the cubic spline of \(P\), with tails and thricing, at \eqref{eq:h-d2}. The pricing error is \(0.035\), \(0.050\), \(1.042\), and \(0.025\). The hold-out error is \(0.039\), \(0.058\), \(1.099\), and \(0.030\). Each of the four curves has one peak on \([0.55F,1.40F]\) and mass \(1.0\). The bandwidth in that row is the same number that prices the chain.

\citeauthor{aitsahalia1998}, at a pricing cross-validation of the implied-volatility regression, has error \(0.064\), \(0.073\), \(1.734\), and \(0.098\), with \(32\), \(14\), \(26\), and \(15\) peaks and mass \(1.6\), \(1.2\), \(1.3\), and \(1.6\). The call cubic of \citeauthor{grith2012} is the smoothest of this group and the closest on price: \(0.031\), \(0.031\), \(0.892\), and \(0.022\), with \(3\), \(3\), \(13\), and \(1\) peaks. Their implied-volatility cubic prices at \(0.030\), \(0.027\), \(0.874\), and \(0.023\), and the integral of the positive part on the window is \(4.3\), \(1.1\), \(1.0\), and \(3.5\).

A peak count treats a ripple and a spike as the same event, and the count depends on where the height cutoff is set. The TV column of Table~\ref{tab:listed} is the total variation of \(\max(\hat q,0)\) on the strike window of Figure~\ref{fig:listed}, divided by twice the highest value on that window. A unimodal curve climbs once and descends once, so the ratio equals one, apart from whatever mass remains at the two edges. An extra oscillation of height \(a\) adds \(2a\) to the numerator and so adds \(a\) divided by the height of the mode. The three bold rows score \(1.0\). The call cubic scores \(1.2\), \(1.1\), \(2.1\), and \(1.2\): a small wiggle on three chains and a material one on NDX. \citeauthor{aitsahalia1998} scores \(10.9\), \(5.0\), \(11.0\), and \(7.6\). \citeauthor{aitsahalia2003} already separate the two widths. The local cubic in Figure~\ref{fig:listed} uses \(0.9\,F\hat\sigma\sqrt{T}\,n^{-1/9}\), about two and a half times \eqref{eq:h-d2}. The call cubic is on the left of Figure~\ref{fig:listed}. Figure~\ref{fig:listed-kern} draws the other two pricing-bandwidth kernels on that same vertical scale.

Averaging the out-of-the-money hold-out error across the four chains, one obtains \(0.306\) for the estimator of this paper, \(1.245\) for \citeauthor{aitsahalia2003}, \(0.774\) for \citeauthor{aitsahalia1998}, \(0.312\) for the call cubic of \citeauthor{grith2012}, \(0.282\) for their implied-volatility cubic, and \(0.336\) for PCA.

\begin{figure}[H]
\centering
\includegraphics[width=0.98\textwidth]{figures/ccdf_listed.pdf}
\caption{SPX, NDX, and RUT, 18 December 2026, snapshot 23~September 2026. Each curve is fit to the projected call of Section~\ref{sec:data}. Left: \(\max(\hat f_{\Q},0)\) for the cubic-spline estimator at \eqref{eq:h-d2}, the lognormal mixture, the \citeauthor{aitsahalia2003} local cubic at \(0.9\,F\hat\sigma\sqrt{T}\,n^{-1/9}\), and the call cubic of \citeauthor{grith2012} at its pricing bandwidth. Right: out-of-the-money implied volatility for the spline, the mixture, the call cubic, and \citeauthor{aitsahalia1998}. The dashed line is the forward. The pricing-bandwidth kernels are in Figure~\ref{fig:listed-kern}.}
\label{fig:listed}
\end{figure}

\begin{figure}[H]
\centering
\includegraphics[width=0.98\textwidth]{figures/ccdf_listed_kern.pdf}
\caption{The pricing-bandwidth kernels on the chains of Figure~\ref{fig:listed}. \citeauthor{aitsahalia1998} and the implied-volatility cubic of \citeauthor{grith2012} use the even/odd pricing bandwidth. Left: \(\max(\hat f_{\Q},0)\), on the strike window and vertical scale of the left panel of Figure~\ref{fig:listed}, with curves truncated there. Right: out-of-the-money implied volatility of the same two fits. The dashed line is the forward.}
\label{fig:listed-kern}
\end{figure}

\section{Conclusion}
\label{sec:conc}

The estimator is a kernel smooth of one expiry at the second-derivative bandwidth \eqref{eq:h-d2}. \citet{aitsahalia1998}, \citet{aitsahalia2003}, \citet{grith2012}, \citet{jackwerth2004}, and \citet{figlewski2018} warn that a kernel density on one expiry is unimodal only after the bandwidth is widened further, and that the bandwidth which prices leaves the peaks. That pricing criterion, applied to the Nadaraya--Watson smile of \citeauthor{aitsahalia1998} and to both local cubics of \citeauthor{grith2012}, reprices as closely or more closely, and the densities have many peaks or an integral well above one on the window past the quotes. The call cubic is the nearest of that group, with \(3\), \(3\), \(13\), and \(1\) peaks. On the four listed chains the spline has one peak and mass \(1.0\), and its out-of-the-money error is below the mixture on each chain. PCA remains a sieve of fixed dimension.

\section*{Code and Data Availability}
A Python implementation of \eqref{eq:thrice} and a replication script are available at \url{https://github.com/s-broda/RND}.

\bibliographystyle{elsarticle-harv}
\bibliography{refs}
\end{document}
